%!TEX program = xelatex
\documentclass[11pt,a4paper]{article}

% -------- Packages --------
\usepackage[margin=0.7in]{geometry}
\usepackage{parskip}
\usepackage{hyperref}
\usepackage{fontspec}
\usepackage{xcolor}
\usepackage{titlesec}
\usepackage{enumitem}

% -------- Font --------
\setmainfont{Times New Roman}

% -------- Colors --------
\definecolor{linkblue}{HTML}{2563eb}
\hypersetup{colorlinks=true, urlcolor=linkblue}

% -------- Formatting --------
\setlist{nosep, left=0pt}
\titlespacing{\section}{0pt}{12pt}{4pt}
\titleformat{\section}{\large\bfseries\uppercase}{}{0pt}{}[\vspace{-2pt}\rule{\textwidth}{0.5pt}]

\begin{document}

% ======== HEADER ========
\begin{center}
{\LARGE\bfseries Ziyu Zhou}
\vspace{4pt}

\small
tzuyu030401@gmail.com $\cdot$ +86 18891550460 $\cdot$ Ningbo, China
\end{center}

% ======== EDUCATION ========
\section{Education}

\textbf{University of Nottingham Ningbo China} \hfill Ningbo, China
\par
\textit{BSc Hons Computer Science with Artificial Intelligence} \hfill Sept 2021 -- July 2026
\begin{itemize}
  \item Full English instruction throughout entire undergraduate study
  \item Awarding body: University of Nottingham, UK
\end{itemize}

% ======== EXPERIENCE ========
\section{Experience}

\textbf{Research Assistant -- Machine Learning Course} \hfill [Semester]
\par\textit{University of Nottingham Ningbo China}
\begin{itemize}
  \item Assisted in preparing tutorial materials, grading assignments, and supporting student learning in Machine Learning coursework
  \item Reinforced understanding of ML algorithms, model evaluation, and data quality --- directly applicable to AI training workflows
\end{itemize}

% ======== LANGUAGE SKILLS ========
\section{Language Skills}

\begin{tabular}{l l}
\textbf{Mandarin Chinese} & Native speaker, standard Northern Mandarin \\
\textbf{Sichuan Dialect (Southwestern Mandarin)} & Native speaker, spoken in Hanzhong (home region) \\
\textbf{Shaanxi Dialect (Central Plains Mandarin)} & Fluent, family language (father from Xi'an) \\
\textbf{English} & IELTS 6.0 (B2) $\cdot$ 4 years of full-English academic environment \\
\end{tabular}

% ======== RELEVANT SKILLS ========
\section{Relevant Skills for AI Tutor}

\begin{itemize}
  \item \textbf{Cross-dialect awareness}: Naturally distinguish between Mandarin, Southwestern Mandarin (Sichuan), and Central Plains Mandarin (Shaanxi) at phonetic, lexical, and prosodic levels
  \item \textbf{Phonetic observation}: Identify tonal variations, stress patterns, intonation contours, and accent features across Chinese dialects and English
  \item \textbf{Transcription}: Accurate character-level transcription across varying audio quality and accent conditions
  \item \textbf{Audio assessment}: Evaluate voice quality, recording environment, background noise, and suitability for model training
  \item \textbf{AI training literacy}: Understand the full pipeline -- data annotation, supervised fine-tuning (SFT), RLHF, and model evaluation -- from CS+AI degree
  \item \textbf{Technical tools}: Familiar with audio annotation platforms, Python, data processing workflows
\end{itemize}

% ======== PORTFOLIO ========
\section{Portfolio Highlights}

\href{https://github.com/zhouziyu12/nlp-chinese}{github.com/zhouziyu12/nlp-chinese}

Prepared a comprehensive annotation portfolio including:
\begin{itemize}
  \item \textbf{Voice samples}: Mandarin, Sichuan dialect, Shaanxi dialect, and English -- demonstrating clear natural delivery across four language varieties
  \item \textbf{Annotated transcripts}: Each sample includes word-level transcription, dialect identification, prosodic analysis, and audio quality assessment
  \item \textbf{Dialect comparison}: Same sentence read in four varieties, with systematic comparison of tonal, lexical, and prosodic differences
  \item \textbf{Practice annotations}: TTS speech quality evaluation, demonstrating independent judgment on synthetic vs. human voice data
\end{itemize}

% ======== RELEVANT EXPERIENCE ========
\section{Academic \& Project Experience}

\textbf{Computer Science with AI Degree Program} \hfill 2021 -- 2026
\par\textit{University of Nottingham Ningbo China}
\begin{itemize}
  \item Coursework in: Machine Learning, Natural Language Processing, Data Structures \& Algorithms, Software Engineering
  \item Developed understanding of AI model training lifecycle, data quality requirements, and evaluation methodologies
  \item All coursework, exams, and presentations conducted in English
\end{itemize}

% ======== ADDITIONAL ========
\section{Additional Information}

\begin{itemize}
  \item \textbf{Time availability}: Immediately available upon graduation (July 2026), open to full-time, part-time, or contract arrangements
  \item \textbf{Equipment}: Personal computer (Windows 10+), stable internet connection, smartphone
  \item \textbf{Location}: Ningbo, China (willing to work remotely globally)
  \item \textbf{No visa sponsorship required} for positions outside the US
\end{itemize}

\end{document}
